%------------------------------------------------------------------------------------
%  SECTION 6.2 : EN-TÊTES ET PIEDS DE PAGE MODULAIRES AVEC SCRLAYER-SCRPAGE
%------------------------------------------------------------------------------------
\section{Styles d'En-têtes et Pieds de Page Modulaires}
\label{sec:styles_en_tetes_pieds}

La gestion des hauts et bas de page en \LaTeX{} a longtemps été confiée au paquet obsolète \texttt{fancyhdr}.
Dans un document moderne basé sur KOMA-Script, \texttt{fancyhdr} provoque des conflits majeurs avec la géométrie et le placement des notes de marge.
\texttt{ClasseGuillaume} s'appuie sur le système de calques haute fidélité \textsf{scrlayer-scrpage}, offrant une modularité sans précédent.

\subsection{Le style par défaut : \texttt{scrheadings}}
\label{subsec:style_scrheadings}

Le style actif par défaut dans l'ensemble du mémoire est \texttt{scrheadings} :
\begin{itemize}
  \item Un filet horizontal fin (\qty{0.5}{pt}) sépare l'en-tête du texte courant.
  \item L'en-tête intérieur affiche dynamiquement le titre du chapitre ou de la section courante via la macro \texttt{\textbackslash headmark}.
  \item L'en-tête extérieur affiche le numéro de page en chiffres gras (\texttt{\textbackslash pagemark}).
  \item Le pied de page reste volontairement épuré pour ne pas alourdir la composition.
\end{itemize}

\subsection{Démonstration en direct du style épuré : \texttt{StyleFocus}}
\label{subsec:style_focus}

Pour les lectures concentrées, les synthèses de fin de partie ou les rapports managériaux, la classe fournit le style \texttt{StyleFocus}.
Ce style supprime le filet d'en-tête, centre le rappel de section entre tirets (\texttt{-- \textbackslash headmark{} --}) et positionne le numéro de page au centre du pied de page.

\begingroup
% Bascule locale en StyleFocus
\pagestyle{StyleFocus}

Nous activons immédiatement \texttt{StyleFocus} sur cette portion de texte.
Observez l'en-tête et le pied de cette page :
la ligne supérieure a disparu, laissant place à une respiration minimale, tandis que la pagination est descendue au centre du pied de page.
Cette disposition est particulièrement prisée pour les sections de réflexion ou les déclarations éthiques.

\subsection{Démonstration en direct du style de diagnostic : \texttt{StyleTest}}
\label{subsec:style_test}

Durant la phase de développement et de révision technique, l'auteur a besoin d'informations de débogage précises : quel est le nom du fichier compilé ? Combien de pages totales compte le document ?

La classe intègre à cet effet le style technique \texttt{StyleTest} :
\begin{itemize}
  \item Un badge orange centré \texttt{\tiny [DEBUG VIEW]} dans l'en-tête.
  \item Des filets supérieur (\qty{0.8}{pt}) et inférieur (\qty{0.4}{pt}).
  \item L'indication explicite du fichier racine (\texttt{\textbackslash jobname.tex}) dans le pied de page.
  \item Le décompte complet des pages via le paquet \textsf{lastpage} : \texttt{p.~\textbackslash thepage/\textbackslash pageref\{LastPage\}}.
\end{itemize}

% Bascule locale en StyleTest
\pagestyle{StyleTest}

Cette nouvelle bascule active instantanément la vue de diagnostic sur les pages courantes.
Elle permet aux équipes de relecture de vérifier en un coup d'œil l'horodatage et la pagination relative sans devoir inspecter les métadonnées internes du fichier PDF.

\endgroup
% Rétablissement automatique de scrheadings à la sortie du groupe
\pagestyle{scrheadings}

La sortie du bloc rétablit automatiquement le style standard \texttt{scrheadings} pour la suite du mémoire.
Ces bascules illustrent la flexibilité extrême offerte par la classe pour rythmer la lecture selon la nature des contenus.

